\documentclass[a4paper,twoside]{article}

\usepackage{morewrites}

\usepackage[svgnames,dvipsnames,x11names]{xcolor}
\usepackage{graphicx}
\usepackage[english]{babel}
\usepackage[colorlinks=true, linkcolor=NavyBlue, urlcolor=NavyBlue, citecolor=NavyBlue]{hyperref}
\usepackage[a4paper]{geometry}
\usepackage{ts-fonts}
\usepackage{booktabs}
\usepackage{csquotes}
\usepackage{tabularx}
\usepackage{tcolorbox}
\usepackage{xcolor}
\usepackage{fvextra}
\usepackage{amsmath}
\usepackage{float}
\usepackage{seqsplit}
\newcolumntype{Y}{>{\centering\arraybackslash}X}
\newcolumntype{Z}{>{\raggedleft\arraybackslash}X}
\usepackage{ts-callouts}
\usepackage{ts-code}

\usepackage{lastpage}
\usepackage{refcount}
\makeatletter
\def\ps@plain{
\let\@oddhead\@empty
\let\@evenhead\@empty
\def\@oddfoot{
\hfil
\ifnum\getpagerefnumber{LastPage}>1\relax
\thepage
\fi
\hfil}
\let\@evenfoot\@oddfoot
}
\makeatother

\title{Emoji and icons}
\author{}\date{}
\begin{document}
\maketitle
Two colon-delimited shortcodes look alike and behave differently, because one
names a \textbf{character} and the other names an \textbf{SVG that only a website has}.
\section{Emoji}
\tscodeinline{:smile:} is sugar for the character it names: the parser replaces it with the
character, and that is what reaches every backend.
\begin{tscode}[lang=md, engine=pygments]
\begin{Verbatim}[commandchars=\\\{\},breaklines, breakanywhere, commandchars=\\\{\}]
Ship it :rocket: :smile: at 12:30:45.
\end{Verbatim}
\end{tscode}
renders as
\begin{displayquote}
Ship it \tsemoji{🚀} \tsemoji{😄} at 12:30:45.
\end{displayquote}
The name table is GitHub's --- the \tscodeinline{gemoji} short names, which \tscodeinline{pymdownx.emoji}
ships as one of its indexes. A colon-delimited word that is \textbf{not} in the
table is literal text with no diagnostic, which is what keeps \tscodeinline{12:30:45} and
\tscodeinline{a:b:c} safe. Nothing fires inside code.

Class E, and GitHub renders the shortcodes too, so the source stays readable
wherever it is pasted. How an emoji is \emph{set} in print --- a colour font, a
monochrome one, an image --- is the template's business, not syntax.
\section{Material icons}
An icon shortcode names an SVG of the MkDocs Material theme, which no print
backend and no icon-less site can honour. Four set prefixes are recognised:
\begin{tscode}[lang=md, engine=pygments]
\begin{Verbatim}[commandchars=\\\{\},breaklines, breakanywhere, commandchars=\\\{\}]
Click :material\PYZhy{}cog: Settings, :fontawesome\PYZhy{}solid\PYZhy{}check: to confirm,
:octicons\PYZhy{}tag\PYZhy{}16: for a release, :simple\PYZhy{}github: for the repository.
\end{Verbatim}
\end{tscode}
Such a shortcode lowers to a \textbf{web-only span} holding the shortcode as its
text. The \tscodeinline{media=web} rule then removes it from print --- the spaces around it
collapse, as for any \href{extra.md\#three-universal-attributes}{zero-width node} ---
and the HTML writer emits \tscodeinline{<span class="icon">:material-\allowbreak{}cog:</span>} for the
site's stylesheet or plugin to replace.

The shortcode is its own canonical spelling: the printer writes such a span
back as the shortcode, so the round-trip is exact.
\begin{tscallout}[kind=warning, title={Icons do not reach the PDF}]
\tscodeinline{tmark check} raises the hint \tscodeinline{icon-\allowbreak{}web-\allowbreak{}only} once per shortcode, so an
author writing for print knows the icon is not there. Icons are decoration;
a symbol that must reach print is an \textbf{emoji} or an \textbf{image}.
\end{tscallout}
\begin{center}
\begin{tabularx}{\linewidth}{>{\raggedright\arraybackslash}X>{\raggedright\arraybackslash}X>{\raggedright\arraybackslash}X>{\raggedright\arraybackslash}X>{\raggedright\arraybackslash}X}
\toprule
\textbf{Construct} & \textbf{Canonical} & \textbf{Print} & \textbf{Web} & \textbf{Class} \\
\midrule
\tscodeinline{:smile:} & the character & the character & the character & E \\
\tscodeinline{:material-\allowbreak{}cog:} & the shortcode & dropped & \tscodeinline{<span class="icon">} & D \\
\bottomrule
\end{tabularx}
\end{center}
See also Smart symbols and quotes for the substitutions that run
on ASCII punctuation.
\end{document}
